\chapter{Conclusion}

We solved the twodimensional system of equations that describe incompressible fluid dynamics.
To do this, we first dealt with the simpler case of linear advection and were able to show, that our methods were accurate comapred to analytical solutions.
Here, we also implemented Lie-Splitting to reduce computational cost.
We then extended our numerical procedure to respect the incompressible characteristics:
Chorins projection method allowed us to reduce the dimensionality of the problem by accepting a numerical error.
Then, pressure was calulated by solving the Poisson equation that could be derived from the incompressibility constraint.
We could show, that the error remained bounded. 
To further minimize it, we implemented spectral methods to solve the Poisson equation.
The resulting divergence could be reduced up to machine precision.
Lastly, these methods turned out to be suitable to reproduce Rayleigh-B\'{e}nard convection cells.
This was done by a certain paramterization and initialization of the model. 
Initial conditions with small perturbations were introduce to trigger the emergence of convective behaviour.
The parameters $\beta= 50.0$ (related to buoyancy) and $\nu = 0.005$ (viscosity) were yielding the desired result and are consistent with Rayleigh theory.
That is why in total, we consider the project to be a successful numerical solution of the incompressible Navier-Stokes problem.


